\documentclass[a4paper, 11pt]{article}
\usepackage{xeCJK}
\usepackage{geometry}
\usepackage{titlesec}
\usepackage{graphicx}   % Required for images
\usepackage{subcaption}
\usepackage[backend=biber, style=numeric]{biblatex}
\addbibresource{../../ref.bib}

\geometry{margin=1in}
\setCJKmainfont{MS Mincho}

\titlespacing*{\section}{0pt}{2ex}{1ex}
\titlespacing*{\subsection}{0pt}{1ex}{0.8ex}
\titlespacing*{\subsubsection}{0pt}{1ex}{0.5ex}

\begin{document}

\begin{center}
    {\LARGE \bfseries 進捗報告} \\ % Title: Large and bold
    \vspace{0.3em} % A small vertical space
    Juan Pablo Corredor \\ % Author
    ２０２５年１２月１２日 % Date
\end{center}
\vspace{0.5em} % Adjust space between title block and main text

\section*{近況}
% Write your recent findings here.
\begin{itemize}
    \item 風邪をひいて、なかなか治らない。
    \item 隅田川都立高等学校との国際交流コラボ。
\end{itemize} 

\section*{やったこと}
\begin{itemize} 
    \item ウェーブレット変換からの低次元表現はまだ開発中。
\end{itemize}

\section*{考察}

Spectral Fluxはやはりトランジェントの初サンプルを見つけるには効果的であるが、トランジェントの終わりを識別するには弱いため、周期認識アルゴリズムの相互相関関数分析で認識できた。

そのために、相互相関のウィンドウ、また音源の波形もゼロ位相フィルタで前処理し、相互相関関数の精度を増やす。

次はトランジェントの分析方法、低次元表現の生成方法等を算段。現在、ウェーブレット変換（CWT）でスケ―ログラム（ウェーブレット変換のスペクトルグラム）を使用し、トランジェントの代表的な特徴が得られる。

\section*{やること}
% Outline future steps here.    
\begin{itemize}
    \item ウェーブレット変換とそのスケ―ログラムを開発。
    \item 上手くいけば、モデル開発に入る。
\end{itemize}

%開放弦

% Add your bibliography here.
\printbibliography

\end{document}